\section{微分形式}
\label{sec:mcalc-forms}

协向量$\alpha\in T_p^*M$把一个切向量读成一个数，它是曲线积分的被积对象：沿曲线累计$\alpha(\dot\gamma)$就得到功或环流。要在曲面上累计通量、在区域上累计体积，被积对象就得读入两个、三个乃至$n$个切向量，读出它们张成的小平行体的“有向大小”。本节先在单个切空间中建立这种交替读数及其乘法，再让它随点光滑变化，最后定义它们沿映射的拉回与外微分。这里只用到光滑结构，不需要度量；函数的微分$df$沿用第~\ref{chap:manifold-basics}章的定义。

\subsection{交替形式与楔积}
\label{sec:mcalc-alternating-wedge}

平面上两个向量$u,v$张成的平行四边形，带符号的面积是$\det[u,v]=u^1v^2-u^2v^1$。这个量对每个向量线性，交换$u,v$时变号，$u,v$共线时为零。符号记录了方向：从$u$逆时针转到$v$为正，顺时针为负。把这三条性质抽象出来，就得到交替形式。

\begin{definition}[交替形式与微分形式]
\label{def:mcalc-alternating-form}
设$V$是$n$维实线性空间，$k\geq1$。$V$上的$k$次\term{交替形式}（alternating form）是多线性映射$\omega:V^k\to\mathbb R$，并且交换任意两个变量时变号：
\[
\omega(\ldots,v_i,\ldots,v_j,\ldots)=-\omega(\ldots,v_j,\ldots,v_i,\ldots).
\]
$k$次交替形式全体记为$\Lambda^k(V^*)$，约定$\Lambda^0(V^*)=\mathbb R$。光滑流形$M$上的$k$次\term{微分形式}（differential form），简称$k$形式，是对每点$p$指定$\omega_p\in\Lambda^k(T_p^*M)$，且在每张坐标图中的系数（见下文）都是光滑函数。$k$形式全体记为$\Omega^k(M)$，$0$形式就是光滑函数。
\end{definition}

交替性有两个直接推论。若两个变量相同，交换它们既不改变$\omega$的值又使其变号，所以值为零；由多线性，只要$v_1,\ldots,v_k$线性相关，$\omega(v_1,\ldots,v_k)=0$。因此$k>n$时$\Lambda^k(V^*)=\{0\}$：多于$n$个向量总是线性相关的。一次交替形式就是线性泛函，$\Lambda^1(V^*)=V^*$。

\paragraph{楔积与坐标基}
从几个线性泛函可以拼出交替形式。对$\alpha^1,\ldots,\alpha^k\in V^*$，定义它们的\term{楔积}（wedge product）
\begin{equation}
(\alpha^1\wedge\cdots\wedge\alpha^k)(v_1,\ldots,v_k)=\det\bigl[\alpha^a(v_b)\bigr]_{a,b=1}^k .
\label{eq:mcalc-wedge-determinant}
\end{equation}
行列式对每列线性、交换两列变号，所以右端是$k$次交替形式。在流形的坐标图中取对偶基$dx^1,\ldots,dx^n$，即$dx^i(\partial_j)=\delta^i_j$，则$dx^1\wedge dx^2(u,v)=u^1v^2-u^2v^1$，正是开头的带符号面积。

\begin{theorem}[交替形式的基]
\label{thm:mcalc-alternating-basis}
设$e_1,\ldots,e_n$是$V$的基，$e^1,\ldots,e^n$是对偶基。对严格递增的指标组$I=(i_1<\cdots<i_k)$记$e^I=e^{i_1}\wedge\cdots\wedge e^{i_k}$。则$\{e^I\}$是$\Lambda^k(V^*)$的基，并且每个$\omega\in\Lambda^k(V^*)$都写成
\[
\omega=\sum_{I}\omega_I\,e^I,\qquad\omega_I=\omega(e_{i_1},\ldots,e_{i_k}).
\]
特别地，$\dim\Lambda^k(V^*)=\binom nk$。
\end{theorem}

\begin{proof}
对递增的$J=(j_1<\cdots<j_k)$，由式\eqref{eq:mcalc-wedge-determinant}，$e^I(e_{j_1},\ldots,e_{j_k})=\det[e^{i_a}(e_{j_b})]$，当$I=J$时为单位矩阵的行列式$1$，当$I\neq J$时某一行全为零，值为$0$。所以$\eta=\sum_I\omega_Ie^I$与$\omega$在每组$(e_{j_1},\ldots,e_{j_k})$上取值相同。两个交替多线性映射在所有递增的基向量组上相同，由交替性在所有基向量组上相同，再由多线性处处相同，故$\omega=\eta$。线性无关性由同一计算得到：若$\sum_I c_Ie^I=0$，在$(e_{j_1},\ldots,e_{j_k})$上取值得$c_J=0$。
\end{proof}

于是在坐标图中，$k$形式写成$\omega=\sum_I\omega_I\,dx^{i_1}\wedge\cdots\wedge dx^{i_k}$，定义中的“系数光滑”指这些$\omega_I$光滑。例如$\mathbb R^3$上的$2$形式有三个系数：$\omega=a\,dy\wedge dz+b\,dz\wedge dx+c\,dx\wedge dy$。

\paragraph{一般楔积及其规则}
式\eqref{eq:mcalc-wedge-determinant}只把一次形式相乘。对$\omega\in\Lambda^k(V^*)$、$\eta\in\Lambda^l(V^*)$，楔积$\omega\wedge\eta\in\Lambda^{k+l}(V^*)$由双线性和下式确定：
\[
(e^{I})\wedge(e^{J})=e^{i_1}\wedge\cdots\wedge e^{i_k}\wedge e^{j_1}\wedge\cdots\wedge e^{j_l}.
\]
等价地，
\begin{equation*}
\begin{aligned}
&(\omega\wedge\eta)(v_1,\ldots,v_{k+l})\\
&\qquad=\frac1{k!\,l!}\sum_\sigma\operatorname{sgn}(\sigma)\,\omega(v_{\sigma(1)},\ldots,v_{\sigma(k)})\,\eta(v_{\sigma(k+1)},\ldots,v_{\sigma(k+l)}),
\end{aligned}
\end{equation*}
其中$\sigma$遍历$\{1,\ldots,k+l\}$的全部置换；两种写法一致，这里不展开核对。微分形式逐点作楔积。由行列式的性质可得下面的规则：
\begin{enumerate}
\item 楔积是双线性的，并满足结合律；
\item 次数交换规则：$\omega\wedge\eta=(-1)^{kl}\,\eta\wedge\omega$；
\item 对一次形式$\alpha$，$\alpha\wedge\alpha=0$。
\end{enumerate}
第二条可以在基上核对：把$e^{j_1},\ldots,e^{j_l}$逐个移过$k$个一次因子，每移过一个变一次号，共变号$kl$次。第三条是第二条在$k=l=1$时的推论。

\begin{example}[楔积的计算]
\label{ex:mcalc-wedge-computation}
在平面上，$(dx+dy)\wedge(dx-dy)=dx\wedge dx-dx\wedge dy+dy\wedge dx-dy\wedge dy=-2\,dx\wedge dy$。在两个向量上取值核对：取$u=\bm e_1$、$v=\bm e_2$，左端为$\det\begin{bmatrix}1&1\\1&-1\end{bmatrix}=-2$，右端为$-2$。
\end{example}

\paragraph{最高次形式与换元}
$n$维流形上的$n$形式只有一个系数：在坐标图中$\omega=f\,dx^1\wedge\cdots\wedge dx^n$。换用坐标$y$时，$dx^i=\sum_a\frac{\partial x^i}{\partial y^a}dy^a$，代入并用规则(2)(3)展开，只有$a$两两不同的项保留，重排后得到
\begin{equation}
dx^1\wedge\cdots\wedge dx^n=\det\Bigl[\frac{\partial x^i}{\partial y^a}\Bigr]\,dy^1\wedge\cdots\wedge dy^n .
\label{eq:mcalc-top-form-change}
\end{equation}
这一步就是行列式的Leibniz展开。系数按Jacobian行列式变换，并且带着符号。下一节的积分正是围绕这个符号展开的：普通的多元积分换元公式使用的是行列式的绝对值。

\subsection{微分形式的拉回}
\label{sec:mcalc-pullback}

第~\ref{chap:manifold-maps}章对光滑映射$F:M\to N$定义了协向量的拉回：$N$上点$F(p)$处的协向量$\alpha$先读$dF_pv$，就得到$p$处的协向量。多个切向量的情形完全相同。

\begin{definition}[微分形式的拉回]
\label{def:mcalc-pullback}
设$F:M\to N$光滑，$\omega\in\Omega^k(N)$。定义$F^*\omega\in\Omega^k(M)$为
\[
(F^*\omega)_p(v_1,\ldots,v_k)=\omega_{F(p)}\bigl(dF_pv_1,\ldots,dF_pv_k\bigr),\qquad v_1,\ldots,v_k\in T_pM .
\]
对$0$形式，$F^*f=f\circ F$。
\end{definition}

由定义直接得到以下规则：$F^*$是线性的；$F^*(\omega\wedge\eta)=F^*\omega\wedge F^*\eta$，因为式\eqref{eq:mcalc-wedge-determinant}中每个$\alpha^a(dF_pv_b)$恰是$(F^*\alpha^a)(v_b)$；对复合映射$G\circ F$，由链式法则$d(G\circ F)_p=dG_{F(p)}\circ dF_p$得$(G\circ F)^*=F^*\circ G^*$，次序与复合相反。$k=1$时它就是协向量的拉回，并满足$F^*(df)=d(f\circ F)$。

\paragraph{坐标计算}
这些规则给出拉回的实际算法：在$N$的坐标$y$中写出$\omega$，把$y^j$换成$F^j(x)$，把$dy^j$换成$d(F^j)=\sum_i\frac{\partial F^j}{\partial x^i}dx^i$，再用楔积规则展开。

\begin{example}[极坐标与圆周上的拉回]
\label{ex:mcalc-pullback-examples}
取$F(r,\theta)=(r\cos\theta,r\sin\theta)$。由$F^*dx=\cos\theta\,dr-r\sin\theta\,d\theta$、$F^*dy=\sin\theta\,dr+r\cos\theta\,d\theta$，
\[
F^*(dx\wedge dy)=r\cos^2\theta\,dr\wedge d\theta-r\sin^2\theta\,d\theta\wedge dr=r\,dr\wedge d\theta ,
\]
其中用了$dr\wedge dr=d\theta\wedge d\theta=0$与$d\theta\wedge dr=-dr\wedge d\theta$。这与式\eqref{eq:mcalc-top-form-change}一致：极坐标变换的Jacobian行列式为$r$。

再取$c:\mathbb R\to\mathbb R^2$，$c(t)=(\cos t,\sin t)$，以及$1$形式$\omega=x\,dy-y\,dx$。代入得$c^*\omega=\cos t\cdot\cos t\,dt-\sin t\cdot(-\sin t)\,dt=dt$。这里$c$不是单射，也不是满射，拉回照样有定义。
\end{example}

\paragraph{为什么沿映射反向作用}
切向量沿$F$从$M$推到$N$，形式则从$N$拉回到$M$。反向并非约定，而是由能否定义决定的。拉回只需要在$p$处读取$dF_p$，对任何光滑映射都有定义；若想把$M$上的形式“推”到$N$，在一点$y\in N$须先选一个原像，而原像可能不存在（$F$不满），也可能有多个且给出不同的值（$F$不单）。例~\ref{ex:mcalc-pullback-examples}中的$c$把$t$与$t+2\pi$送到同一点，$\mathbb R$上的形式无法据此唯一地送到圆周上。这与第~\ref{chap:manifold-maps}章中向量场的推送需要额外条件是同一个原因。

\subsection{外微分}
\label{sec:mcalc-exterior-derivative}

函数的微分$df$把$0$形式变成$1$形式，记录函数沿每个方向的一阶变化。我们希望对任意$k$形式也有这样的运算，把它变成$k+1$形式，使Stokes公式中“内部的导数”有意义。在坐标中，最直接的做法是对系数求微分，再与原有的基形式作楔积。

\begin{definition}[外微分]
\label{def:mcalc-exterior-derivative}
在坐标图中，对$\omega=\sum_I\omega_I\,dx^{i_1}\wedge\cdots\wedge dx^{i_k}$定义
\begin{equation}
d\omega=\sum_Id\omega_I\wedge dx^{i_1}\wedge\cdots\wedge dx^{i_k}
=\sum_I\sum_{j=1}^n\frac{\partial\omega_I}{\partial x^j}\,dx^j\wedge dx^{i_1}\wedge\cdots\wedge dx^{i_k}.
\label{eq:mcalc-exterior-derivative}
\end{equation}
$d:\Omega^k\to\Omega^{k+1}$称为\term{外微分}（exterior derivative）。
\end{definition}

对$0$形式$f$，公式给出$df=\sum_j\partial_jf\,dx^j$，正是函数的微分，所以这里没有重新定义它。平面上$d(P\,dx+Q\,dy)=\partial_yP\,dy\wedge dx+\partial_xQ\,dx\wedge dy=(\partial_xQ-\partial_yP)\,dx\wedge dy$，括号恰是Green公式中出现的组合。$\mathbb R^3$上，若把$1$形式$F^1dx+F^2dy+F^3dz$与$2$形式$F^1dy\wedge dz+F^2dz\wedge dx+F^3dx\wedge dy$都与向量场$(F^1,F^2,F^3)$对应，则对函数、$1$形式、$2$形式取$d$，系数分别是梯度、旋度和散度。

\begin{theorem}[外微分的性质]
\label{thm:mcalc-exterior-derivative}
式\eqref{eq:mcalc-exterior-derivative}定义的$d$与坐标的选取无关，因而给出整体的线性映射$d:\Omega^k(M)\to\Omega^{k+1}(M)$，并满足：
\begin{enumerate}
\item 对函数，$d$就是函数的微分；
\item 乘积规则：$d(\omega\wedge\eta)=d\omega\wedge\eta+(-1)^k\omega\wedge d\eta$，其中$\omega$为$k$形式；
\item $d\circ d=0$；
\item 与拉回相容：对光滑$F:M\to N$，$F^*(d\omega)=d(F^*\omega)$。
\end{enumerate}
\end{theorem}

\begin{proof}
先在单张坐标图中验证(2)(3)。对(2)，由线性只需考虑$\omega=f\,dx^I$、$\eta=h\,dx^J$。$\omega\wedge\eta=fh\,dx^I\wedge dx^J$，$d(fh)=h\,df+f\,dh$，于是$d(\omega\wedge\eta)=h\,df\wedge dx^I\wedge dx^J+f\,dh\wedge dx^I\wedge dx^J$。第二项中把$dh$移过$k$个一次因子，产生符号$(-1)^k$，得到$(-1)^k\omega\wedge d\eta$；第一项就是$d\omega\wedge\eta$。

对(3)，先看函数：$d(df)=\sum_{i,j}\partial_i\partial_jf\,dx^i\wedge dx^j$。二阶偏导数关于$i,j$对称，而$dx^i\wedge dx^j$关于$i,j$反对称，成对抵消，所以$d(df)=0$。由定义，坐标基形式$dx^I$的系数为常数，$d(dx^I)=0$；对一般的$\omega=f\,dx^I$，由(2)，$d(d\omega)=d(df\wedge dx^I)=d(df)\wedge dx^I-df\wedge d(dx^I)=0$。

再证(4)在坐标中成立。由线性和拉回保持楔积，只需检查$\omega=f\,dy^{j_1}\wedge\cdots\wedge dy^{j_k}$。$F^*\omega=(f\circ F)\,d(F^{j_1})\wedge\cdots\wedge d(F^{j_k})$，对其取$d$，由(2)和(3)，各个$d(F^{j_a})$再求外微分为零，只剩$d(f\circ F)\wedge d(F^{j_1})\wedge\cdots$。而$d(f\circ F)=F^*df$，所以结果是$F^*(df\wedge dy^{j_1}\wedge\cdots)=F^*(d\omega)$。

最后，坐标变换是两组坐标之间的微分同胚，同一个形式在两组坐标中的表达互为拉回。由(4)，先变换再求$d$与先求$d$再变换结果相同，所以两张图中由式\eqref{eq:mcalc-exterior-derivative}定义的$d\omega$在重叠处一致，拼成整体的形式。整体上的(1)—(4)由局部成立推出。
\end{proof}

\paragraph{闭形式与恰当形式}
满足$d\omega=0$的形式称为\term{闭形式}（closed form），可以写成$\omega=d\eta$的称为\term{恰当形式}（exact form）。由$d\circ d=0$，恰当形式都是闭的。反过来未必成立，障碍来自流形的拓扑。在穿孔平面上取
\[
\omega=\frac{x\,dy-y\,dx}{x^2+y^2}.
\]
直接计算可得$d\omega=0$。若$\omega=df$，则沿单位圆$c(t)=(\cos t,\sin t)$有$c^*\omega=c^*df=d(f\circ c)$；由例~\ref{ex:mcalc-pullback-examples}，在单位圆上$c^*\omega=dt$，于是$\frac{\mathrm d}{\mathrm dt}f(c(t))=1$，$f(c(2\pi))-f(c(0))=2\pi$，与$c(2\pi)=c(0)$矛盾。$\omega$在每个不含整条射线的区域上是辐角函数的微分，但辐角无法在整个穿孔平面上单值地定义。闭形式与恰当形式之差刻画的“洞”，与第~\ref{chap:homology}章用同调刻画的是同一类信息，这里不展开二者的对应。

\paragraph{与切映射的区别}
外微分$d$与映射的微分$dF_p$都写作$d$，作用对象却不同。$dF_p$由点映射$F:M\to N$得到，作用于切向量，输出$N$的切向量；外微分作用于$M$上的一个形式场，输出次数加一的形式场，不涉及第二个流形。二者唯一的交汇是$0$形式：实值函数$f$既可以看成映射$f:M\to\mathbb R$，又可以看成$0$形式，此时$df_p$作为切映射与$df$作为外微分给出同一个线性读数。
